\section{Zadanie 16}
Wskaż, jak można wykorzystać eksperyment, w którym rzucamy \(n\)
razy igłą długości 1 na płaszczyznę z równoległymi liniami narysowanymi
o odstępach jeden,
do oszacowania wartości \(\frac{1}{\pi}\). Jak duże powinno być \(n\),
aby z prawdopodobieństwem \(0.95\) nasze oszacowanie było z~dokładnością
do \(0.01\)? Skorzystaj z CTG.

\subsection*{Rozwiązanie}
Nasze zdarzenie to standardowa igła Buffona,
czyli każdy rzut jest próbą Bernoulliego o prawdopobieństwie
sukcesu \(\frac{2}{\pi}\).

Wartość oczekiwana sumy tych indykatorów po n próbach jest
\(E(\sum_{i=1}^n X_i) = \frac{2}{\pi}n\),\\
a wariancja wyniesie
\(Var(\sum_{i=1}^n X_i) = n \cdot \frac{2}{\pi} \cdot
\left(1-\frac{2}{\pi}\right)\)

Dla średniej mamy: \(E(S) = \frac{2}{\pi}\),
\(Var(S) = \frac{1}{n} \cdot \frac{2}{\pi} \cdot
\left(1-\frac{2}{\pi}\right)\)

Zamiast obliczać \(\frac{1}{\pi}\) z dokładnością \(0.01\)
obliczymy \(\frac{2}{\pi}\) z dokładnością \(0.02\).

Przyjmiemy \(p=\frac{2}{\pi}\), \(\varepsilon = 0.02\)

\[
	P(p-\varepsilon \leqslant S \leqslant p+\varepsilon)
	= P(-\varepsilon \leqslant S-p \leqslant \varepsilon)
\]
\[G
	= P\left(
	-\varepsilon \cdot \sqrt{\frac{\pi^2n}{2(\pi - 2)}}
	\leqslant (S-p) \cdot \sqrt{\frac{\pi^2n}{2(\pi - 2)}}
	\leqslant \varepsilon \cdot \sqrt{\frac{\pi^2n}{2(\pi - 2)}}
	\right)
\]

Kopiując z zadania wyżej (lub sprawdzając w tablicach),\\ widzimy że
\(\Phi(x) \geq 0.975\) dla \(x \geq 1.959964\). Czyli
\[
	0.02 \cdot \sqrt{\frac{\pi^2n}{2(\pi - 2)}}
	\geqslant 1.959964 \implies n \geqslant 2222
\]

Moglibyśmy założyć, że mamy jakąś ograniczoną wiedzę na
temat \(\pi\) (gdybyśmy znali dokładną wartość \(\pi\), obliczenie
\(\frac{1}{\pi}\) mogłoby być wykonane prościej).

Wiedząc jedynie, że \(\pi > 2\), dostalibyśmy \(n \geqslant 2401\)
(pesymistycznym przypadkiem jest tutaj \(\pi = 4\))



\subsection*{Odpowiedź}
\[
	\boxed{\text{2222}}
\]
